\histitem{谐波分析与群论的联系}

我们在《积分学》一书的历史注记中（参见 ÉHM，第275页）描述过，与热方程有关的问题如何引导傅里叶提出把任意函数表示为三角函数之和这一革命性思想。

我们无法在此追溯整个十九世纪傅里叶级数和积分的发展，也无法追溯该理论提出的问题对分析演变的深刻影响——其影响甚至颠覆了实数的观念，并促成集合论的诞生（参见 ÉHM，第42页）。但为了理解本书所阐述思想的形式，有必要说明谐波分析在1925年以后如何同群论和希尔伯特谱理论发生联系。

傅里叶思想与 R/Z 和 R 的群结构之间联系的认识，以及将这些思想推广到其他群，来得相当晚。

在有限群及其特征标理论的早期阶段，自然可以辨认出一些如今看来包含有限傅里叶分析雏形的结果。例如，Z/qZ 的可逆元群的对偶，以及表达一个元素的特征函数为特征标线性组合的正交关系，已隐约出现在1837年狄利克雷的文章 {[}17{]} 中；该文证明当 a 与 q 互素时，存在无穷多个素数 \(p \equiv a \mod q\) 。他借助高斯的思想；高斯在研究具有整数系数及固定判别式的二元二次型时引入了“特征标”一词（参见 {[}22，第230节{]}）。

但高斯和狄利克雷都没有使用群的语言。戴德金在1879年讲述这些工作时，才认识到它们潜在的作用，并为交换群定义了特征标的概念。H. 韦伯大大扩充了戴德金的评述：先是1882年 {[}77{]}，然后是1886年 {[}78，第112页{]}；他在那里引入有限阿贝尔群 A 的对偶群，并注意到后者同构于 A。他当时的目的是构造具有给定伽罗瓦群的域 Q 的阿贝尔扩张，并未考虑这类群上的谐波分析问题。

在非交换群情形，有限群表示理论的基础在20世纪初前后，继 G. 弗罗贝尼乌斯、W. 伯恩赛德和 I. 舒尔的工作之后得到坚实奠定（参见 ÉHM，第154页）。1905年以后，特征标的正交关系在理论组织中发挥关键作用已很清楚。但它们与谐波分析的亲缘关系仍隐藏不显。

我们应归功于赫尔曼·外尔将这些领域联系起来：他构想了紧群表示的一般理论。在1925年至1927年间发表的几部重要著作中，他认识到这类理论与傅里叶分析及希尔伯特谱理论之间的联系，并为许多后来的发现奠定基础。

是舒尔的一封信把外尔引向这条道路的，缘于二人对不变量理论的共同兴趣。我们在关于 Haar 测度的注记（ÉHM，第289–291页）中描述过，A. 胡尔维茨早在1897年就已构想利用“不变积分”来确定在正交群下不变的 \(R^{n}\) 上多项式。舒尔似乎直到1924年才知道胡尔维茨的结果；那时他突然理解到，不变积分可用于把自己为有限群建立的特征标理论和正交关系推广到群 \(\mathbf{O}(n)\)。外尔立即看出，舒尔的方法与埃利·嘉当的工作相结合，使得可以构造紧半单李群的不可约表示。外尔将嘉当和舒尔的思想融为一体的三篇论文发表于1925年 {[}83{]}，其中已经包含 LIE, IX, §6–7 的基本结果；参见 ÉHM，第328–330页。

但在次年一篇与学生 F. 彼得合写且同样著名的文章 {[}54{]} 中，外尔才揭示紧群表示论、傅里叶分析和希尔伯特谱理论之间的联系。这篇不足二十页的文章孕育了许多未来的发展。

彼得和外尔摆脱了对所研究群的结构所作的任何代数假设，只假定 G 是一个配备不变测度的紧拓扑群。对连通李群而言，这种测度的存在是显然的；不久之后 A. Haar 将在一般情形建立它，部分动机正来自此处呈现的结果（参见 ÉHM，第291页）。

他们研究的出发点是矩阵系数的正交性。舒尔注意到，对每个不可约表示等价类选取一个代表 \(\pi \in \widehat{\mathrm{G}}\) 、\(\pi\) 的空间 E 上一个 G-不变内积以及 E 的一个标准正交基，就得到一族在 \(\mathrm{L}^2 (\mathrm{G})\) 中正交归一的矩阵系数。彼得和外尔计划证明这族函数是完备的。

在有限群情形，弗罗贝尼乌斯曾通过研究群代数 C{[}G{]} 证明了相应结果。彼得和外尔接着考虑 G 上连续函数空间 \(\mathcal{C}(\mathrm{G})\) ，并为其赋予卷积乘法。这似乎是首次将卷积乘法作为抽象运算、明确地同群结构联系起来（参见 ÉHM，第295页）。此外，彼得和外尔把 \(\mathcal{C}(\mathrm{G})\) 的一个元素称为 Gruppenzahl（“群数”），并以 xy 表示两个 Gruppenzahlen 的（卷积）乘积。他们还为 \(\mathcal{C}(\mathrm{G})\) 配备对合 \(f \mapsto \widetilde{f}\) ，其中 \(\widetilde{f}(g) = \overline{f(g^{-1})}\) ，对于 \(g \in G\) 。

在引入对合代数 \(\mathcal{C}(\mathrm{G})\) 后，彼得和外尔的第一个观察是，每个 G 的酉表示 \(\pi\) 都给出一个表示 \(f \mapsto \pi(f)\) ，其代数为 \(\mathcal{C}(\mathrm{G})\) （参见 V，第400页）。他们明确给出了 \(\mathcal{C}(\mathrm{G})\) 的埃尔米特元素概念，也（隐含地）给出了正元素概念，为应用希尔伯特空间技巧打开道路。

彼得和外尔不再犹豫，称算子 \(\pi(f)\) 为 \(f\) 的“傅里叶系数”，并通过指出舒尔正交关系蕴含贝塞尔不等式，继续与傅里叶理论作平行比较

\[
\sum_ {\pi \in \widehat {\mathrm{G}}} \dim (\pi) \operatorname{Tr} (\pi (f) \pi (f) ^ {*}) \leqslant (f * \widetilde {f}) (e) = \| f \| _ {2} ^ {2}.\tag{6}
\]

希尔伯特--施密特谱理论随后使他们能够证明这是等式（“普朗谢雷尔公式”）。对 \(\varphi\) （\(\mathcal{C}(\mathrm{G})\) 中的每个元素），彼得和外尔关联连续核 \(K_{\varphi}\colon G\times G\to C\) ，其定义为 \(\mathrm{K}_{\varphi}(x,y)=\varphi(xy^{-1})\) 。随后他们观察到，如果 \(\varphi\) 形如 \(f*\widetilde{f}\) 且 \(f\in\mathcal{C}(\mathrm{G})\) ， \(^{(6)}\) 则 \(K_{\varphi}\) 是希尔伯特--施密特理论意义下的正埃尔米特核。施密特算法的一个变体 \(^{(7)}\) 使他们可以构造核 \(K_{\varphi}\) 的特征函数标准正交基，继而将(6)中等式的证明归结为：由 \(K_{\varphi}\) 定义的算子的迹等于其特征值之和。

当然，只有表示 \(\pi\) 满足 \(\pi(f) \neq 0\)才可能出现在(6)中；\(K_{\varphi}\) 的谱理论表明，只要 \(f\) 不恒为零，就存在这样的表示。彼得和外尔很容易由此推出，不可约表示分离 G 的点——这是盖尔范德--赖科夫定理的一个特例，稍后将讨论该定理。

在弗罗贝尼乌斯的著名工作中，与(6)类似的等式曾使人能够证明 G 的正则表示包含所有不可约表示。但这涉及将该等式用于 C{[}G{]} 的单位元 f；此时显然有 \(\pi(f) \neq 0\) 对所有 \(\pi\) 都成立 。依彼得和外尔之见，卷积没有单位元解释了他们证明中的许多困难；然而，他们借用了傅里叶级数理论中一个注定具有巨大影响的思想，指出可将(6)应用于一列构成卷积逼近单位的函数，从而推出弗罗贝尼乌斯结果的类似物（参见 I，第120页第10号）。

同样地，从(6)产生的等式出发，其极化形式

\[
\sum_ {\pi \in \widehat {\mathrm{G}}} \dim (\pi) \operatorname{Tr} (\pi (x) \pi (y)) = (x * y) (e)\tag{7}
\]

通过令 \(y\) 为卷积逼近单位，可得到 \(\mathcal{C}(\mathrm{G})\) 的每个元素都能由不可约表示的矩阵系数的线性组合一致逼近。这些方法刚刚使外尔能对 H. 玻尔的“几乎周期函数”理论的基本结果给出新证明 {[}84{]}；我们很快将回到该理论。彼得和外尔指出，其中可以看出对非紧群表示理论的第一个分析应用，这里该群是 \(\mathbf{R}\) 的平移群。

年轻的量子理论很快给外尔提供了重新回到 R 的表示的机会。1927年底，他注意到（无论有限与否的）群和表示对于阐明该理论基础的意义 {[}85{]}。次年，他在一本轰动性的著作 {[}86{]} 中再次讨论这一问题。

在1927年的文章中，外尔注意到，若 u 是希尔伯特空间上的连续自伴算子，则 \(t \mapsto e^{itu}\) 是 R 的酉表示。但存在并非这种形式的 R 的酉表示；重新采用冯·诺伊曼关于用对称部分算子表示可观测物理量的思想，外尔提出每个 R 的酉表示都形如 \(t \mapsto e^{itu}\) ，其中 u 是希尔伯特空间上的部分自伴算子。这一观点使他能够把海森堡和薛定谔的“正则关系”（冯·诺伊曼已作详细研究）的分析，归结为同一希尔伯特空间上两个 R 的酉表示之间的联系。斯通 {[}73{]} 在1932年、沿着他对自伴算子谱性质的研究，证明了外尔所期待的结果（V，第428页，定理1）。

